


\documentclass{article}
\usepackage{enumerate}
\usepackage[margin=1.4in, top=1in]{geometry}
\title{Teaching Statement}
\author{Joel Villatoro}
\date{}

\usepackage{blindtext}
\usepackage{hyperref}
\begin{document}
\vspace{-0.5in}
\maketitle

\section*{Introduction}
As an academic mathematician studying pure mathematics, I consider teaching and mentoring as a core component of my job. 
The core of my teaching philosophy can be described by my two main focuses, student engagement and accessibility. 
By student engagement, I mean getting students involved in the course beyond a lecture/homework format. Accessibility means remaining available and approachable to my students and mentees.

My experiences teaching cover a wide range of ages, skill levels, backgrounds, and formats. However, I believe that I still have much to learn and great teaching and mentorship is not a specific goal but a lifelong pursuit. At the University of Illinois I was successful as a TA and was awarded for excellence. At KU Leuven in Belgium I was faced with a significantly different student body and classroom environment and that presented me with many new challenges. At Washington University, due to restrictions placed by my NSF grant, I was limited in the amount of paid teaching I could do but I did teach a large lecture course as well as an upper division course with a smaller class size. In addition to classroom work, at all of these institutions I have also engaged with undergraduates in a mentor/advising context. I will first discuss my views and philosophy regarding classroom instruction and then I will discuss my experiences as a mentor.

\section*{Classroom methodology}


My views on how to maintain an engaging classroom have been shaped by my experiences. My formative teaching experience was as a teaching assistant at the University of Illinois. At UIUC, the main approach was active learning, which meant that the students were split into small groups and given a worksheet. My role was as a facilitator for discussions and provide students with hints and suggestions if they get stuck. In this time I realized the importance of learning to identify struggling or distracted students and to take proactive action rather than simply waiting for questions. In my experience, this format was generally very successful and, at first, it was the only teaching format I was familiar with.

Unsurprisingly, when it came time for me to work as the primary instructor, I found that maintaining student engagement was significantly more challenging. Thus far,
I have employed a few strategies to break up the tedium of lectures and improve student engagement. My favorite and generally most successful method is the creation of interactive visualizations of mathematical concepts. For example, while teaching differential equations, I created an interactive phase diagram visualizer\footnote{link: \url{https://www.desmos.com/calculator/tj3vn0rcqj}]} so that students could develop intuition for integral curves and eigenvectors. Another example is python notebook that illustrates a topological proof of the fundamental theorem of algebra \footnote{link:\url{https://colab.research.google.com/drive/1A898TEOGSHCNcYKM0BWP8YmsE2PLFz9M?usp=sharing}} which I created for another instructor. 
I am also frequently on the lookout for online resources (such as videos or other multimedia) that can compliment the lectures and worksheets. Especially for undergraduate level topics, there is a vast quantity of high quality visualizations and instructional resources. I intend to continue to experiment with and learn new techniques to provide an engaging student-centered classroom.

In any mathematics course, the majority of learning will inevitably take place by doing exercises and peer discussion. For courses with large enrollment, this means that the exercise sections and teaching assistants have a particularly significant impact on student outcomes. When I taught differential equations the students reported that they felt they learned the most from doing the worksheets I wrote and from interacting with their TA. While this is not surprising, it highlights the fact that primary instructors for such classes have their work cut out for them when it comes to directly impacting student success. Many of the levers they have available to them are fundamentally indirect such as preparing course materials or supervising and training TAs.

When teaching a class, I try to frequently remind students that they are welcome to come to my office and email me with their concerns.
However, I have found that this is not always enough as some students require more proactive measures to ensure that they are getting the help they need. After each exam, I try to identify which students appear to be struggling and I send them a personal invitation to visit me to discuss the class. It's important to approach these interactions with sensitivity because students can often feel embarrassed or demoralized when they don't achieve the success they envisioned for themselves. For example, I avoid using phrases such as "you appear to be struggling in the class" as it can be seen as overly personal and assigning blame to the student. I also try to remain cognizant of and sensitive to the diverse backgrounds of my students.
This means taking advantage of common ground I may have with other students as well as showing respect for their cultural and personal experiences of my students.

\section*{Mentorship}
I would like to now take the time to discuss some of my views on and experiences in mentoring/advising students. Broadly, my philosophical takeaways from these experiences is that, when mentoring any student, it is important to take into account and focus on the student's goals for themselves. I have also found that, when working with undergraduates, one of the best ways to connect with my mentees is to learn alongside them.

At UIUC, I worked a mentor for undergraduates through the Illinois Geometry Lab. The Illinois Geometry Lab is a program run by the mathematics department which pairs small groups of undergraduates with a faculty member and graduate student. The faculty member provides the undergraduates with a project, usually a visualization of some mathematical concept, while the graduate student provides regular support while the undergraduates work on the project. I worked on multiple projects during my time but the first one involved numerical simulations of differential equations. At the time, I had little background in programming or computational mathematics. Despite this, I found that this was in many ways a benefit as I was better able to bond and connect with the undergraduate students by learning along side them. 

At KU Leuven my primary work as a mentor was supervising a Bachelor project. This involved advising a group of undergraduate mathematics majors while they studied chosen topic in mathematics. This experience was distinct from the IGL as the students were mathematics majors primarily interested in pursuing a career in pure mathematics. Much of our activities were reading and discussing mathematics articles and excerpts from textbooks. My primary role was to guide and advise them as they dipped their toe into reading and learning about topics in more advanced mathematics. The specific topic we investigated was cystallography and Bieberbach groups which was not a topic I was directly familiar with. Once again, I was reminded of how rewarding it can be to learn alongside my students.

Overall, I feel that my past mentorship experiences have provided me with the tools to be able to flexibly work with undergraduate students in a way that takes into account their personal interests and objectives.

\section*{Diversity Equity and Inclusion}

I have always seen promoting equity and diversity as a vital component of my work. I am deeply committed to promoting an environment that is welcoming to individuals of all backgrounds. At your university, I will commit to furthering its mission to provide an equitable and inclusive environment for its faculty and students.

As a discipline, mathematics is faced with an ongoing crisis of diversity and equity. To be able to address this problem, it is vital to acknowledge the struggles faced by women and members of underrepresented ethnic and cultural groups. I recognize that people of many backgrounds often face serious challenges such as a lack of role models, implicit bias, and individual and systemic discrimination. 

There are a few different ways I try to promote inclusiveness and equity when teaching and mentoring students. One way is by recognizing how certain class formats can disproportionately advantage or disadvantage students based on their backgrounds. Generally speaking, traditional lectures result in worse outcomes for URM students. On the other hand, active learning has been shown to narrow achievement gaps for underrepresented students in mathematics\footnote{See \url{https://doi.org/10.1073/pnas.1916903117}}. 

In order to reap the benefits of more equitable formats, it is necessary for the instructor to be committed to providing an inclusive classroom environment. One way I try to accomplish this is by making sure that I talk to every student, rather than favoring certain students. I also try to show an active interest in my students' backgrounds and life experiences. I believe that students are much more likely to feel as if they belong if they have established a rapport with their colleagues. Therefore, I try to encourage student interactions both inside and outside the classroom. For example, I encourage students to form study groups and to work together when doing homework exercises. When splitting students into groups in class, I prefer to assign groups randomly to avoid the formation of cliques which can make students from underrepresented groups feel excluded.

In mathematics, I believe that maintaining a ``growth mindset'' is particularly important to establishing an inclusive classroom environment. To elaborate, women and students from certain underrepresented groups sometimes struggle with an internalized ``fixed mindset'' which tells them that their ability to succeed in mathematics is intrinsic and unchangeable. Although all students can experience these feelings, I have found that they are particularly pernicious among students from underrepresented groups. For example, I have seen this frequently with the high school students I have interacted with through the Puertas program. Many of them, especially young Latina women, have sadly internalized a belief that they are simply ``not good at math'' and that this is an immutable characteristic. It is important to confront this attitude head on and encourage students to instead adopt a growth mindset where they view mathematical ability as a skill that can be improved.

\section*{Conclusion}

In conclusion, my teaching and mentoring philosophy centers on student engagement, accessibility, and inclusivity. I am committed to creating dynamic and interactive classrooms while supporting individual student goals. Additionally, I am dedicated to promoting diversity and equity in mathematics education. I look forward to the opportunity to build upon these skills and make a positive difference for my students.
\end{document}
